\chapter{Verified Simulation Edges}\label{ch:verified}

\section{TM to Generalized Shift (Moore Theorem 7)}

\begin{proposition}[TM $\to$ GS]\label{prop:tm-gs}
\lean{TuringMachineToGeneralizedShift.fromBiTM}
\lean{TuringMachineToGeneralizedShift.stepCommutes}
\lean{TuringMachineToGeneralizedShift.decodeEncode}
\lean{TuringMachineToGeneralizedShift.mooreHaltingForward}
\leanok
\uses{simulation-def, overhead-def, bi-tm, gs-def}
Any $(s,k)$ bi-infinite TM is simulated by a generalized shift with
window width~3, alphabet size $sk+k$, $\sovh = 1$, $\tovh = 1$.

The encoding merges the TM state into the tape cell at the head position:
the extended alphabet has $k$ passive symbols (no state) and $sk$ active
symbols encoding (state, color) pairs. The GS window is
$[\mathrm{left}, \mathrm{active}, \mathrm{right}]$.

Key theorems:
\begin{itemize}
\item \texttt{decodeActiveEncodeActive}: alphabet encoding roundtrip
\item \texttt{decodeEncode}: full configuration roundtrip (nonempty tapes)
\item \texttt{stepCommutes}: one TM step = one GS step
\item \texttt{stepCommutesNorm}: step commutation up to tape normalization
\item \texttt{mooreHaltingForward}: halting preservation
\end{itemize}
Fully verified: zero \texttt{sorry}.
\end{proposition}

\section{GS to Turing Machine (Moore Theorem 8)}

\begin{proposition}[GS $\to$ TM]\label{prop:gs-tm}
\lean{GeneralizedShiftToTuringMachine.toBiTM}
\lean{GeneralizedShiftToTuringMachine.buildTransition}
\lean{GeneralizedShiftToTuringMachine.decodeConfig_encodeConfig}
\notready
\uses{simulation-def, overhead-def, bi-tm, gs-def}
Any GS with window width $w$ and max shift magnitude $m$ is simulated by a
bi-infinite TM with $\sovh = 1$, $\tovh \leq 2(w-1) + m$.

The TM executes in three phases per GS step:
\begin{enumerate}
\item Read phase ($w-1$ steps): scan right, accumulate window contents in
  state encoding.
\item Write phase ($w-1$ steps): write replacement symbols right-to-left.
\item Shift phase ($\leq m$ steps): move head to match GS shift.
\end{enumerate}

Proved: window encoding roundtrip (\texttt{decodeWindow\_encodeWindow}),
configuration roundtrip (\texttt{decodeConfig\_encodeConfig}).
Pending: step simulation for the active case (\texttt{stepSimulation\_w1},
1~\texttt{sorry}).
\end{proposition}

\section{ECA Rule 110 and Rule 124 (Mirror)}

\begin{proposition}[ECA mirror bisimulation]\label{mirror-sim}
\lean{mirrorProperty, mirrorSimulationSteps}
\leanok
\uses{simulation-def, eca-def}
Rule~110 and Rule~124 are related by tape reversal, giving a bisimulation with
$\sovh = 1$, $\tovh = 1$.

The identity $\mathrm{rule110}(a,b,c) = \mathrm{rule124}(c,b,a)$ holds for all
$a,b,c \in \{0,1\}$ (\texttt{mirrorProperty}, by \texttt{fin\_cases} +
\texttt{decide}).
\texttt{reverseTapeInvolutive}: double reversal is identity.
\texttt{mirrorSimulationGeneral}: one step commutes with reversal.
\texttt{mirrorSimulationSteps}: $k$ steps commute with reversal.
Fully verified: zero \texttt{sorry}.
\end{proposition}

\section{2-Tag to Cyclic Tag System (Cook 2004)}

\begin{proposition}[Tag $\to$ CTS]\label{prop:tag-cts}
\lean{tagToCyclicTagSystem}
\lean{tagToCyclicTagSystemSimulation}
\lean{tagToCyclicTagSystemHaltingForward}
\lean{cyclicTagSystemToTagHalting}
\leanok
\uses{simulation-def, tag-def, cts-def}
Any 2-tag system with alphabet size $k$ is simulated by a cyclic tag system
with $2k$ appendants. One tag step corresponds to exactly $2k$ CTS steps.

The encoding maps each symbol $i \in \{0,\ldots,k-1\}$ to a one-hot binary
word of length $k$ (1 at position $i$, 0 elsewhere). The CTS has $k$
appendants encoding the tag productions and $k$ empty appendants.

Key theorems:
\begin{itemize}
\item \texttt{tagToCyclicTagSystemSimulation}: 1 tag step = $2k$ CTS steps
\item \texttt{tagToCyclicTagSystemHaltingForward}: tag halts $\Rightarrow$ CTS halts
\item \texttt{cyclicTagSystemToTagHalting}: CTS halts $\Rightarrow$ tag halts
\end{itemize}
Fully verified in both directions: zero \texttt{sorry}.
\end{proposition}

\section{The Cocke--Minsky Chain}

\begin{proposition}[TM $\to$ Tag $\to$ CTS]\label{prop:cocke-minsky}
\lean{cockeMinskyTag}
\lean{cockeMinskyConfigurationEncode}
\lean{turingMachineToCyclicTagSystem}
\uses{simulation-def, bi-tm, tag-def, cts-def, prop:tag-cts}
\notready
The Cocke--Minsky encoding maps any TM to a 2-tag system with alphabet
size $sk + k + 1$, using markers $A(q,a)$ for state-symbol pairs, $B(a)$
for tape cells, and $C$ as separator.

The composition \texttt{turingMachineToCyclicTagSystem} chains
TM~$\to$~Tag~$\to$~CTS. The halting correspondence
(\texttt{cockeMinskyHaltingForward}) is fully proved, but the step simulation
(\texttt{CockeMinskyStepSimulation}) remains a hypothesis.
\end{proposition}

\section{Main Universality Theorem}

\begin{theorem}[Wolfram's $(2,3)$ TM is universal]\label{thm:wolfram23}
\lean{wolfram23Universal}
\uses{prop:cocke-minsky, wolfram-23-tm}
\notready
For any TM $M$ and initial configuration $c$, there exists a configuration of
Wolfram's $(2,3)$ machine that simulates the computation of $M$ on $c$.

The proof composes:
TM $\xrightarrow{\text{Cocke--Minsky}}$ 2-Tag
$\xrightarrow{\text{Cook}}$ CTS
$\xrightarrow{\text{Smith}}$ TM(2,3).

Currently requires two hypotheses:
\texttt{CockeMinskyStepSimulation} (TM $\to$ Tag step binding) and
\texttt{SmithSimulation} (CTS $\to$ TM(2,3) via Smith's 6-level hierarchy).
\end{theorem}
